\documentclass[10pt]{article}
\usepackage[T1]{fontenc}
\usepackage[utf8]{inputenc}
\usepackage[french]{babel}
\usepackage{amsmath,amssymb,mathrsfs,mathpazo}
\usepackage{geometry}
\geometry{a4paper,margin=2cm}
\usepackage{graphicx,array,tabularx,tabularray,colortbl,multirow,multicol,booktabs,enumitem,pifont,ulem,fancybox,diagbox,cancel}
\usepackage{pstricks,pstricks-add,pst-plot,pst-tree,pst-node,pst-text,pst-func,pst-math,pst-bspline,pst-3dplot,pst-eucl,pst-solides3d}
\let\pscalculate\undefined
\usepackage{pst-bezier}
\usepackage{tikz,pgfplots,tkz-tab}
\usetikzlibrary{arrows,arrows.meta,patterns,calc,positioning,trees,decorations.markings,shapes.geometric,shapes.misc,shapes.arrows,intersections,decorations.pathmorphing,decorations.pathreplacing,backgrounds,shadows,babel}
\pgfplotsset{compat=1.18}
\usepgfplotslibrary{fillbetween}
\usepackage[np]{numprint}
\usepackage{listings,makecell,pas-tableur,twoopt}
\let\barre\undefined
\usepackage[pstricks]{tablvar}
\definecolor{darkgreen}{RGB}{0,100,0}
\providecommand{\red}{\color{red}}
\providecommand{\blue}{\color{blue}}
\providecommand{\green}{\color{green}}
\providecommand{\black}{\color{black}}
\providecommand{\white}{\color{white}}
\providecommand{\gray}{\color{gray}}
\providecommand{\cyan}{\color{cyan}}
\providecommand{\magenta}{\color{magenta}}
\providecommand{\yellow}{\color{yellow}}
\providecommand{\violet}{\color{violet}}
\providecommand{\orange}{\color{orange}}
\providecommand{\euro}{\texteuro}
\providecommand{\EUR}{\texteuro}
\providecommand{\ieme}{\textsuperscript{e}}
\providecommand{\degres}{\ensuremath{^\circ}}
\providecommand{\degree}{\ensuremath{^\circ}}
\providecommand{\vv}[1]{\overrightarrow{#1}}
\providecommand{\Coord}[1]{#1}
\providecommand{\decosix}{}
\providecommand{\starredbullet}{\ensuremath{\bullet}}
\pagestyle{empty}
\setlength{\parindent}{0pt}
\usepackage{adjustbox}\DeclareUnicodeCharacter{2236}{:}\DeclareUnicodeCharacter{2190}{\ensuremath{\leftarrow}}
\providecommand{\R}{\mathbb{R}}
\providecommand{\N}{\mathbb{N}}
\providecommand{\Z}{\mathbb{Z}}
\providecommand{\Q}{\mathbb{Q}}
\providecommand{\D}{\mathbb{D}}
\providecommand{\C}{\mathbb{C}}
\providecommand{\e}{{\mathrm{e}}}
\providecommand{\vect}[1]{\overrightarrow{#1}}
\providecommand{\vectt}[1]{\overrightarrow{\mathrm{#1}}}
\providecommand{\vvt}[1]{\overrightarrow{\mathrm{#1}}}
\providecommand{\np}[2][]{{#2~#1}}
\providecommand{\vv}[1]{{\overrightarrow{#1}}}
\renewcommand{\d}{{\mathrm{d}}}
\providecommand{\degres}{{^\circ}}
\providecommand{\degre}{{^\circ}}
\providecommand{\blue}{\color{blue}}
\providecommand{\red}{\color{red}}
\providecommand{\enskip}{\quad}
\providecommand{\llbracket}{\lbrack\!\lbrack}
\providecommand{\rrbracket}{\rbrack\!\rbrack}
\providecommand{\euro}{{\text{€}}}
\providecommand{\no}{{\text{n°}}}
\providecommand{\interval}[2]{{[#1;#2]}}
\providecommand{\intervFO}[2]{{[#1;#2[}}
\providecommand{\dotfill}{\ldots}

\providecommand{\ds}{\displaystyle}
\providecommand{\textding}[1]{#1}
\providecommand{\Oij}{\(\mathrm{O};\vec{\imath},\vec{\jmath}\)}
\providecommand{\Oijk}{\(\mathrm{O};\vec{\imath},\vec{\jmath},\vec{k}\)}
\providecommand{\Ouv}{\(\mathrm{O};\vec{u},\vec{v}\)}
\providecommand{\intervFF}[2]{\ensuremath{[#1;#2]}}
\renewcommand{\euro}{\ifmmode\text{\texteuro}\else\texteuro\fi}
\DeclareUnicodeCharacter{00B0}{\ifmmode^{\circ}\else\textdegree\fi}
\DeclareUnicodeCharacter{20AC}{\ifmmode\text{\texteuro}\else\texteuro\fi}
% Énoncés : Éducation nationale. Transcription et corrigés : APMEP, Denis Vergès et contributeurs. Découpage et classement : CoopMaths, Sébastien Lozano, Éric Elter et contributeurs. Adaptation de la présentation : Maths Explorer.
% CC BY-SA : https://coopmaths.fr/www/about/
% https://coopmaths.fr/alea/static/eam/2026/tex/eam_specialite_2026_09_antilles_1_cor.tex
\begin{document}
\noindent\textit{APMEP / CoopMaths. CC BY-SA. Présentation adaptée par Maths Explorer.}\par\medskip
\begin{enumerate}
	\item Pour une fonction du second degré $f(x) = ax^2+bx+c$, l'abscisse du sommet de la parabole
	est $x_S = -\dfrac{b}{2a}$. Ici $a = 1$ et $b = 2$, donc
	\[x_S = -\dfrac{2}{2 \times 1} = -1 \qquad\text{et}\qquad f(-1) = (-1)^2 + 2\times(-1) - 3 = 1 - 2 - 3 = -4.\]
	Le sommet de la parabole $\mathcal{P}$ a pour coordonnées $(-1 ; -4)$.

	\item
	\begin{enumerate}
		\item $g(-3) = (-3)^2 + (-3) - 6 = 9 - 3 - 6 = 0$ \quad et \quad $g(2) = 2^2 + 2 - 6 = 4 + 2 - 6 = 0$.

		Les réels $-3$ et $2$ sont donc racines de la fonction $g$.

		\item $g$ est une fonction du second degré de coefficient dominant $1 > 0$ : $g$ est positive à l'extérieur de ses racines et négative entre les
		deux.

		\begin{center}
			\begin{tikzpicture}
				\tkzTabInit[lgt=3.4,espcl=2.6]{$x$/1,$g(x)$/1}{$-\infty$,$-3$,$2$,$+\infty$}
				\tkzTabLine{,+,z,-,z,+,}
			\end{tikzpicture}
		\end{center}
	\end{enumerate}

	\item Pour tout réel $x$ :
	\[f(x) - h(x) = \left(x^2 + 2x - 3\right) - (x+3) = x^2 + 2x - 3 - x - 3 = x^2 + x - 6 = g(x).\]

	\item Étudier la position relative de deux courbes revient à étudier le signe de la différence
	$f(x) - h(x) = g(x)$, étudié à la question \textbf{2. b.}
	
Pour tout $x \in ]-3~;~2[$, $g(x) <0 \iff f(x)-h(x) <0 \iff f(x) < h(x)$.
Donc $\mathcal{P}$ est en dessous de $\mathcal{D}$ sur $]-3~;~2[$.
Et $\mathcal{P}$ est au-dessus de $\mathcal{D}$ sur $]-\infty~;~-3[ \,\cup\, ]2~;~+\infty[$.

	Les points d'intersection correspondent aux annulations de $g$, c'est-à-dire $x = -3$ et
	$x = 2$ : $h(-3) = 0$ et $h(2) = 5$.

Les deux courbes se coupent aux points de coordonnées $(-3~;~0)$ et $(2~;~5)$.
	\item
	\begin{center}
		\begin{tikzpicture}[scale=0.8, x=7mm,y=7mm,>=stealth]
			\draw[step=1,gray!75,line width=0.5pt] (-5,-5) grid (4,6);
			\draw[->,line width=1pt] (-5.3,0) -- (4.3,0);
			\draw[->,line width=1pt] (0,-5.3) -- (0,6.3);
			\foreach \x in {-5,-4,-3,-2,-1,1,2,3,4} \draw (\x,2pt) -- (\x,-2pt) node[below] {\footnotesize \np{\x}};
			\foreach \y in {-5,-4,-3,-2,-1,1,2,3,4,5,6} \draw (2pt,\y) -- (-2pt,\y) node[left] {\footnotesize \np{\y}};
			\draw (-2pt,-2pt) node[below left]{\footnotesize 0};
			\draw[red,line width=1.2pt,domain=-4.15:2.15,samples=80]
			plot(\x,{(\x)^2+2*(\x)-3}) node[left,red]{$\mathcal{P}$};
			\draw[blue,line width=1.2pt] (-5,-2) -- (3,6) node[right,blue]{$\mathcal{D}$};
			\fill (-1,-4) circle(2pt) node[left]{$(-1;-4)$};
			\fill (-3,0) circle(2pt) node[above left]{$(-3~;~0)$};
			\fill (2,5) circle(2pt) node[right]{$(2~;~5)$};
		\end{tikzpicture}
	\end{center}
\end{enumerate}

\end{document}
